% =============================================================================
% Templaté do Relatório Final do Projeto de Aprendizagem de Máquina (pós).
% Cada seção traz comentarios "% TODO grupo:" indicando o que escrever.
% Limite recomendado: 25 paginas (sem apêndices).
% Compilar com: pdflatex relatório.tex (duas passadas)
% =============================================================================
\documentclass[11pt,a4paper]{article}
% Preambulo compartilhado dos artefatos LaTeX da disciplina de Aprendizagem
% de Máquina (graduação e pós-graduação).
%
% Regras editoriais: sem emojis, sem em-dashes, portugues do Brasil com
% acentuacao correta. Arquivos em UTF-8 sem BOM.
%
% Conjunto de pacotes deliberadamente enxuto para garantir compilacao com
% pdflatex em diferentes instalacoes (TeX Live, MiKTeX), sem tikz nem
% microtype.

\usepackage[utf8]{inputenc}
\usepackage[T1]{fontenc}
\usepackage[brazilian]{babel}
\usepackage{geometry}
\geometry{a4paper,margin=2.5cm}
\usepackage{graphicx}
\usepackage{booktabs}
\usepackage{xcolor}
\usepackage{enumitem}
\usepackage{caption}
\usepackage[hidelinks]{hyperref}

% Cores institucionais aproximadas (CIn-UFPE).
\definecolor{cinazul}{RGB}{0, 51, 102}
\definecolor{cinverde}{RGB}{0, 102, 51}

% Macro para titulo de etapa do projeto: Etapa 1 (teórica) e Etapa 2
% (experimental). Substitui o macro \fase do projeto-MDA, mantendo o mesmo
% padrao de marcacao no sumário.
\newcommand{\etapa}[2]{%
    \section*{Etapa #1: #2}%
    \addcontentsline{toc}{section}{Etapa #1: #2}%
}

% Macro para destacar tarefa pendente do grupo em uma caixa visivel.
\newcommand{\tarefagrupo}[1]{%
    \par\medskip
    \noindent\fbox{%
        \begin{minipage}{\dimexpr\textwidth-2\fboxsep-2\fboxrule\relax}
        \textbf{Tarefa do grupo:} #1
        \end{minipage}%
    }%
    \par\medskip
}

% Macro para citar arquivos do repositório em fonte monoespacada.
\newcommand{\arquivo}[1]{\texttt{#1}}

\hypersetup{
    pdfauthor={Prof. Leandro Maciel Almeida},
    pdfkeywords={TabArena, classificação tabular, aprendizagem de máquina}
}


\hypersetup{pdftitle={Relatório Final - Projeto de Aprendizagem de Máquina (pós) - 2026.1}}

\title{\bfseries Relatório Final do Projeto\\[0.3em]
       \large Disciplina de Aprendizagem de Máquina (pós-graduação)}
\author{Grupo: \textbf{[NOMES DOS INTEGRANTES]}\\
        Modelo atribuído: \textbf{[NOME DO MODELO]}}
\date{2026.1}

\begin{document}

\maketitle
\thispagestyle{empty}

\begin{abstract}
% TODO grupo: escreva 5 a 8 linhas resumindo o modelo atribuído, o conjunto
% de datasets do TabArena utilizado, as métricas avaliadas, o resultado
% comparativo contra os baselines (LightGBM, CatBoost, XGBoost) e contra o
% AutoGluon 1.4 (default e extreme), e a principal conclusão sobre os
% regimes em que o modelo se mostra competitivo.
\noindent
[Resumo do projeto: modelo, dados, métricas, comparação com baselines e
AutoGluon, conclusão principal.]
\end{abstract}

\tableofcontents
\bigskip

% =============================================================================
\section{Introdução}
% =============================================================================

\subsection{Contexto e motivação}
% TODO grupo: situe o problema da classificação tabular como dominio em
% que métodos baseados em árvores ainda dominam, mas com novos paradigmas
% emergindo (foundation models, MLPs modernas, redes recursivas, etc.).
% Justifique a relevancia da comparação sistematica.

\subsection{Modelo atribuído}
% TODO grupo: nome, ano, autores, URL do repositório oficial, toolkit
% utilizado e por que o modelo e considerado representativo da literatura
% recente (2023 a 2026).

\subsection{Contribuições do trabalho}
% TODO grupo: liste em 3 a 5 itens o que este relatório entrega
% (implementação reproduzível, comparação com baselines + AutoGluon,
% análise estatística classica e Bayesiana, análise por regime, model
% card detalhado).

% =============================================================================
\etapa{1}{Fundamentação teórica do modelo}
% =============================================================================

\subsection{Motivação e intuição}
% TODO grupo: descreva em 1 a 2 paragrafos a intuição por tras do modelo.
% Qual problema o modelo tenta resolver de forma diferente dos baselines?

\subsection{Arquitetura e formulação}
% TODO grupo: detalhe a arquitetura. Inclua um diagrama de blocos
% (figura) e a formulação matemática essencial. Use \begin{figure}
% \includegraphics{...} \end{figure} para o diagrama.

\subsection{Complexidade computacional}
% TODO grupo: tempo e memória em função de n (amostras) e p (features).
% Liste o pico de memória observado nos seus experimentos (RAM ou VRAM).

\subsection{Posicionamento na literatura tabular SOTA}
% TODO grupo: situe o modelo na linha do tempo 2023 a 2026. Quais são os
% trabalhos correlatos? Quais são as alternativas que o modelo supera ou
% complementa? 5 a 10 referências relevantes.

\subsection{Hiperparâmetros principais e faixas}
% TODO grupo: tabela com hiperparâmetros do modelo, valores tipicos e
% faixas usadas na busca via Optuna.

\begin{table}[h]
\centering
\caption{Hiperparâmetros principais do modelo atribuído e faixas de busca.}
\label{tab:hp}
\begin{tabular}{l p{4.5cm} p{6.0cm}}
\toprule
\textbf{Hiperparâmetro} & \textbf{Tipo / faixa} & \textbf{Justificativa} \\
\midrule
{}[hp1] & [int / log uniform / etc.] & [referência ou heuristica] \\
{}[hp2] & [...]                       & [...] \\
{}[hp3] & [...]                       & [...] \\
\bottomrule
\end{tabular}
\end{table}

% =============================================================================
\etapa{2}{Estudo experimental}
% =============================================================================

\subsection{Métodologia}
% TODO grupo: descreva o protocolo: 30 datasets do TabArena-v0.1
% estratificados em 10 pequenos + 10 médios + 10 grandes; split 70/30
% estratificado por classe com seed fixa documentada; busca de
% hiperparâmetros via Optuna (TPE) com validação cruzada no treino;
% métricas AUC OvO, ACC, G-Mean, Cross-Entropy e tempo total
% (tune + train + predict).

\subsection{Baselines e referência AutoML}
% TODO grupo: descreva os três baselines (LightGBM, CatBoost, XGBoost
% via pytabkit com defaults TD) e as duas configurações do AutoGluon
% 1.4 (preset default e preset extreme com orçamento de 4 horas).

\subsection{Resultados gerais}
% TODO grupo: tabela com media +/- desvio em cada métrica por modelo,
% mais ranking médio nos 30 datasets. Inclua todos os 15 sistemas
% (10 modelos atribuíveis + 3 baselines + 2 AutoGluon) para
% comparação agregada.

\begin{table}[h]
\centering
\caption{Resultados agregados: media +/- desvio padrao em 30 datasets.}
\label{tab:resultados}
\footnotesize
\begin{tabular}{l c c c c c}
\toprule
\textbf{Sistema} & \textbf{AUC OvO} & \textbf{ACC} & \textbf{G-Mean} & \textbf{Cross-Ent.} & \textbf{Tempo (s)} \\
\midrule
{}[Modelo do grupo]    & 0,xx $\pm$ 0,0xx & 0,xx $\pm$ 0,0xx & 0,xx $\pm$ 0,0xx & 0,xx $\pm$ 0,0xx & xxx \\
LightGBM (TD)        & 0,xx $\pm$ 0,0xx & 0,xx $\pm$ 0,0xx & 0,xx $\pm$ 0,0xx & 0,xx $\pm$ 0,0xx & xxx \\
CatBoost (TD)        & 0,xx $\pm$ 0,0xx & 0,xx $\pm$ 0,0xx & 0,xx $\pm$ 0,0xx & 0,xx $\pm$ 0,0xx & xxx \\
XGBoost (TD)         & 0,xx $\pm$ 0,0xx & 0,xx $\pm$ 0,0xx & 0,xx $\pm$ 0,0xx & 0,xx $\pm$ 0,0xx & xxx \\
AutoGluon default    & 0,xx $\pm$ 0,0xx & 0,xx $\pm$ 0,0xx & 0,xx $\pm$ 0,0xx & 0,xx $\pm$ 0,0xx & xxx \\
AutoGluon extreme 4h & 0,xx $\pm$ 0,0xx & 0,xx $\pm$ 0,0xx & 0,xx $\pm$ 0,0xx & 0,xx $\pm$ 0,0xx & xxx \\
{}[demais 9 modelos]   & ... & ... & ... & ... & ... \\
\bottomrule
\end{tabular}
\end{table}

\subsection{Análise por regime}
% TODO grupo: segmente os resultados em quatro eixos: (i) tamanho do
% dataset (n menor que 1k, 1k a 10k, acima de 10k); (ii) número de
% classes (binário vs multiclasse); (iii) proporção categorica
% (baixa vs alta); (iv) presenca de valores ausentes (com vs sem
% NaN). Inclua um gráfico para cada eixo (figura gerada pelo notebook
% \arquivo{notebooks/04\_demo\_stats\_regime.ipynb}).

\subsection{Análise estatística}
% TODO grupo: aplique o protocolo de Demsar (2006) e Benavoli et al.
% (2017): teste de Friedman global, post-hoc Nemenyi (diagrama de
% diferença crítica gerado por \texttt{autorank}), e teste Bayesian
% signed-rank com região de equivalencia prática (ROPE) sobre AUC,
% gerado por \texttt{baycomp}. Discuta os pares com p > 0.95 de
% equivalencia prática.

\subsection{Custo vs desempenho}
% TODO grupo: scatter de tempo total (eixo x) por AUC OvO media (eixo
% y) com todos os 15 sistemas. Discuta a fronteira de Pareto.
% AutoGluon define o teto prático (com orçamento longo); os baselines
% definem o piso barato. Posicione o modelo do grupo nesse mapa.

% =============================================================================
\section{Discussão}
% =============================================================================

% TODO grupo: 4 a 6 paragrafos. Em quais regimes o modelo do grupo se
% mostra competitivo? Em quais perde para os baselines? Em quais perde
% para o AutoGluon? Como interpretar a comparação Bayesiana com ROPE
% nesse contexto? Quais são os limites práticos do modelo?

% =============================================================================
\section{Reprodutibilidade}
% =============================================================================

% TODO grupo: ambiente Python (versão), bibliotecas e versões,
% hardware utilizado, comandos mínimos para reproduzir, hash do
% commit do repositório e tempo total de execução.

\begin{verbatim}
uv sync
python -m src.pipeline.run_all --include-group-model --seed 42
\end{verbatim}

% =============================================================================
\section*{Referências}
% =============================================================================

% TODO grupo: liste artigos, documentacoes e bibliotecas. Use \texttt{bibtex}
% ou liste manualmente.
\begin{itemize}[leftmargin=1.2em]
\item [paper original do modelo]
\item Demsar, J. (2006). Statistical comparisons of classifiers over multiple datasets. JMLR.
\item Benavoli, A., Corani, G., Demsar, J., Zaffalon, M. (2017). Time for a Change: a Tutorial for Comparing Multiple Classifiers Through Bayesian Analysis. JMLR.
\item TabArena-v0.1 (NeurIPS 2025): https://tabarena.ai
\end{itemize}

\appendix

% =============================================================================
\section{Tabela dos 30 datasets escolhidos}
% =============================================================================

% TODO grupo: preencha com os 30 datasets escolhidos do TabArena-v0.1
% (10 pequenos + 10 médios + 10 grandes), citando OpenML task ID, n,
% n_features, n_classes e regime atribuído.

\begin{table}[h]
\centering
\caption{Datasets do TabArena-v0.1 utilizados, agrupados por regime.}
\label{tab:datasets}
\footnotesize
\begin{tabular}{l r r r r l}
\toprule
\textbf{Dataset} & \textbf{Task ID} & \textbf{n} & \textbf{features} & \textbf{classes} & \textbf{Regime} \\
\midrule
{}[d1]  & xxxxx & xxx     & xx & x & pequeno \\
{}[d2]  & xxxxx & xxx     & xx & x & pequeno \\
\multicolumn{6}{c}{...} \\
{}[d10] & xxxxx & xxx     & xx & x & pequeno \\
{}[d11] & xxxxx & xxxxx   & xx & x & médio \\
\multicolumn{6}{c}{...} \\
{}[d20] & xxxxx & xxxxx   & xx & x & médio \\
{}[d21] & xxxxx & xxxxxxx & xx & x & grande \\
\multicolumn{6}{c}{...} \\
{}[d30] & xxxxx & xxxxxxx & xx & x & grande \\
\bottomrule
\end{tabular}
\end{table}

% =============================================================================
\section{Espacos de busca utilizados pelo Optuna}
% =============================================================================

% TODO grupo: para cada modelo cuja busca foi feita pelo grupo, liste o
% espaco de busca efetivamente utilizado (hiperparâmetro, tipo, faixa,
% prior). Para os baselines com defaults TD do pytabkit, basta indicar
% que a configuração default foi usada sem busca.

\end{document}
